% META: {"id": "1SPE-SECDEG-EX-019", "chapitre": "1SPE-SECOND-DEGRE", "type_objet": "exercice", "capacites": ["1SPE-SECOND-DEGRE-C1", "1SPE-SECOND-DEGRE-C2"], "capacites_codes": ["C1", "C2"], "methodes": ["M1", "M2"], "parcours": 2, "competences": ["calculer", "raisonner", "communiquer"], "duree_min": 20, "mode_creation": "ex_nihilo", "sources_inspiration": [], "parametres_sympy": {}, "fichier_tex": "chapitres/1SPE-SECOND-DEGRE/exercices/1SPE-SECDEG-EX-019.tex", "corrige_tex": "chapitres/1SPE-SECOND-DEGRE/corriges/1SPE-SECDEG-CO-019.tex", "status": "generated"}
% BEGIN-VERIFY
% from sympy import symbols, expand, simplify, Rational
% x = symbols('x')
% # f(x) = 3x^2 - 6x - 24
% f = 3*x**2 - 6*x - 24
% # forme canonique : sommet alpha = 6/(2*3) = 1, beta = f(1) = 3 - 6 - 24 = -27
% alpha = Rational(6, 6)
% beta = f.subs(x, alpha)
% assert alpha == 1
% assert beta == -27
% assert simplify(f - (3*(x-1)**2 - 27)) == 0
% # forme factorisee : 3*(x-4)*(x+2), racines 4 et -2
% from sympy import factor
% assert simplify(factor(f) - 3*(x - 4)*(x + 2)) == 0
% # variations : a=3>0, decroissante sur (-inf,1], croissante sur [1,+inf)
% # minimum = -27
% END-VERIFY
\begin{exercice}{1SPE-SECDEG-EX-019}{2}{20}
On considère la fonction $f$ définie sur $\mathbb{R}$ par $f(x) = 3x^2 - 6x - 24$.

\begin{enumerate}
  \item \textbf{(C1)} Identifier les coefficients $a$, $b$, $c$ et écrire $f$ sous forme canonique.
  \item \textbf{(C2)} En déduire les coordonnées du sommet $S$ de la parabole représentative de $f$ et dresser le tableau de variations de $f$ sur $\mathbb{R}$.
  \item \textbf{(C1)} Calculer $\Delta$ et factoriser $f(x)$.
  \item \textbf{(C2)} Déterminer graphiquement (puis algébriquement) les abscisses des points d'intersection de la parabole avec l'axe des abscisses.
  \item L'équation $f(x) = -27$ admet-elle des solutions ? Justifier à l'aide du tableau de variations.
\end{enumerate}
\end{exercice}
